\chapter{Evaluation}

\section{Evaluation Goals}
This chapter evaluates whether the implemented AI collaborator addresses the core objectives formulated in the introduction. Specifically, the evaluation seeks to answer the following research questions:
\begin{itemize}
    \item \textbf{RQ1 (AI Generation \& Efficiency):} To what extent does the integration of an AI-based collaborator reduce the time and manual effort required to construct MILP models compared to purely manual development?
    \item \textbf{RQ2 (Validation \& Accuracy):} How accurately can the AI agent generate mathematically and syntactically valid model structures, and how effectively does the system's validation layer prevent erroneous modifications?
    \item \textbf{RQ3 (Human-in-the-Loop \& Feedback):} How effectively does the iterative human-in-the-loop feedback mechanism enable the operator to recover from AI-generated errors and converge on a correct formulation?
\end{itemize}

To answer these questions, we evaluate the system architecture using a benchmark dataset of optimization problems tested across multiple Large Language Models (LLMs) of varying sizes.

\section{Experimental Setup}
\subsection{System Configuration}
All experiments are conducted on the integrated prototype described in this thesis, using the identical execution pipeline for all tasks. To understand how the system's performance scales with the underlying intelligence of the agent, the evaluation is performed across three different LLMs accessed via OpenRouter, representing small, medium, and large model classes. Each run uses the same system-prompt contract and identical tool availability to ensure configuration consistency.

The runtime environment consists of the modeler web client session, a LangGraph-backed orchestration runtime for threaded inference \cite{langchainLangGraph}, and a browser-automation service using Puppeteer for model read/write operations \cite{pptrPuppeteerPuppeteer}.

\subsection{Dataset and Task Design}
Instead of relying on a user study, the evaluation employs a dataset-driven benchmarking approach. The dataset consists of 10 representative MILP modeling tasks:
\begin{itemize}
    \item \textbf{7 Easy Problems:} Standard formulations (e.g., simple Knapsack, Diet Problem) requiring fewer variables and constraints.
    \item \textbf{3 Hard Problems:} Complex formulations (e.g., Facility Location, Multi-commodity Flow) requiring advanced structural adjustments and multiple interacting constraints.
\end{itemize}

For each problem in the dataset, two execution modes are measured:
\begin{enumerate}
    \item \textbf{Manual mode:} The operator models the problem entirely through direct UI interaction.
    \item \textbf{AI-assisted mode:} The operator models the problem with the AI collaborator through natural-language requests, providing iterative feedback when necessary.
\end{enumerate}

\subsection{Evaluation Procedure}
For each sample problem, the operator attempts to generate a correct, feasible model using both modes. During the AI-assisted mode, the initial prompt is provided, and the system's first generated output is evaluated. If the model is incorrect or syntactically invalid, the operator provides corrective natural-language feedback until a valid model is achieved. 

Throughout this process, the following data points are recorded: manual modeling time, AI modeling time, first-go correctness, validation layer interventions, and the number of feedback iterations required.

\section{Evaluation Metrics}

\subsection{Efficiency Metrics (RQ1)}
\begin{itemize}
    \item \textbf{Manual Modeling Time:} Total time required to manually build the model.
    \item \textbf{AI-Assisted Modeling Time:} Total time required to build the model using the AI, including prompt writing and review time.
    \item \textbf{Speedup Factor:} The ratio of manual time to AI-assisted time.
\end{itemize}

\subsection{Accuracy and Validation Metrics (RQ2)}
\begin{itemize}
    \item \textbf{First-Go Correctness:} The percentage of tasks where the AI generated a mathematically correct and syntactically valid model on its first attempt.
    \item \textbf{Validation Catch Rate:} The number of times the system's validation layer successfully intercepted and blocked an invalid or hallucinated artifact before it could corrupt the workspace.
\end{itemize}

\subsection{Human-in-the-Loop Metrics (RQ3)}
\begin{itemize}
    \item \textbf{Correction Iterations:} For models not correct on the first go, the average number of additional human-feedback prompts required to fix the errors.
    \item \textbf{Final Success Rate:} The percentage of tasks successfully completed after iterative human refinement.
\end{itemize}

\section{Results}
\subsection{RQ1: AI Generation \& Efficiency}
To answer RQ1, we compare the manual and AI-assisted completion times across the easy and hard datasets for the different LLMs.

\begin{table}[h]
\centering
\caption{Efficiency comparison across model sizes and task complexities.}
\label{tab:rq1_efficiency}
\begin{tabular}{lcccc}
\hline
\textbf{Model Size} & \textbf{Task Level} & \textbf{Manual Time (avg)} & \textbf{AI Time (avg)} & \textbf{Speedup} \\
\hline
Small LLM & Easy (7) & -- & -- & -- \\
Small LLM & Hard (3) & -- & -- & -- \\
Medium LLM & Easy (7) & -- & -- & -- \\
Medium LLM & Hard (3) & -- & -- & -- \\
Large LLM & Easy (7) & -- & -- & -- \\
Large LLM & Hard (3) & -- & -- & -- \\
\hline
\end{tabular}
\end{table}

Table~\ref{tab:rq1_efficiency} presents the core efficiency metrics recorded during the evaluation. The \textit{Manual Time} serves as the baseline, representing the operator's speed when using the platform without AI assistance. The \textit{AI Time} captures the end-to-end duration from submitting the initial natural language prompt to the final acceptance of the model, including any necessary review or wait times. The \textit{Speedup} factor is calculated as the ratio of manual time to AI time. A speedup factor greater than 1.0 indicates that the AI collaborator successfully accelerated the modeling process. By categorizing the results into Easy and Hard tasks, this table allows us to observe whether the efficiency gains are consistent across varying levels of problem complexity, and whether larger, slower LLMs negatively impact overall modeling speed due to high inference latency.

\subsection{RQ2: Validation \& Accuracy}
To answer RQ2, we measure the system's ability to generate valid models on the first attempt and the effectiveness of the validation layer in intercepting errors.

\begin{table}[h]
\centering
\caption{First-go accuracy and validation performance across LLMs.}
\label{tab:rq2_accuracy}
\begin{tabular}{lccc}
\hline
\textbf{Model Size} & \textbf{First-Go Correct (\%)} & \textbf{Caught by Val.} & \textbf{Uncaught Errors} \\
\hline
Small LLM & -- & -- & -- \\
Medium LLM & -- & -- & -- \\
Large LLM & -- & -- & -- \\
\hline
\end{tabular}
\end{table}

The metrics in Table~\ref{tab:rq2_accuracy} illustrate the raw generation capability and safety of the architecture. \textit{First-Go Correct (\%)} indicates the proportion of tasks where the AI successfully generated a fully valid and mathematically accurate model without requiring any human intervention. When the model failed to produce a correct formulation on the first attempt, the errors were either intercepted by the system or reached the user. \textit{Caught by Val.} tracks the number of times the architectural validation layer successfully identified and blocked syntactically invalid or hallucinated artifacts, proving the effectiveness of the system's safety mechanisms. Conversely, \textit{Uncaught Errors} represents logical flaws that bypassed technical validation and required human identification, directly informing the necessity of the human-in-the-loop workflow evaluated in the next section.

\subsection{RQ3: Human-in-the-Loop \& Feedback}
To answer RQ3, we analyze the correction overhead required when the AI fails to produce a correct model initially.

\begin{table}[h]
\centering
\caption{Iterative correction cycles required to achieve a valid model.}
\label{tab:rq3_hitl}
\begin{tabular}{lccc}
\hline
\textbf{Model Size} & \textbf{Avg. Corrections} & \textbf{Max Iterations} & \textbf{Final Success (\%)} \\
\hline
Small LLM & -- & -- & -- \\
Medium LLM & -- & -- & -- \\
Large LLM & -- & -- & -- \\
\hline
\end{tabular}
\end{table}

Table~\ref{tab:rq3_hitl} evaluates the effectiveness of the iterative feedback mechanism for the subset of tasks that were not solved correctly on the first attempt. The \textit{Avg. Corrections} metric tracks the average number of follow-up prompts the operator had to provide to guide the AI toward a correct solution. A low average indicates that the AI agent successfully maintains context and incorporates human feedback efficiently. \textit{Max Iterations} highlights the worst-case scenario encountered during testing, indicating potential failure states or "prompt loops." Finally, the \textit{Final Success (\%)} demonstrates the ultimate reliability of the platform; it shows the percentage of tasks that were eventually solved correctly when combining the AI's generation capabilities with human supervision, regardless of initial failures.

\section{Threats to Validity}
\subsection{Internal Validity}
The evaluation relies on a predefined dataset of optimization problems. The prompts used to describe these problems may contain implicit biases based on the operator's formulation style. To mitigate this, standard problem descriptions from optimization literature are used where possible.

\subsection{External Validity}
The 10 selected problems represent a subset of possible MILP formulations and may not generalize to highly specific industrial applications. Furthermore, \textbf{LLM Bias} acts as a significant threat to external validity: the architecture's output quality is inherently bounded by the capabilities of the underlying LLMs. While testing across three model sizes mitigates this by showing performance variance, future advancements or regressions in LLM capabilities will directly impact the measured efficiency and accuracy of the platform.

\subsection{Construct Validity}
Measuring efficiency via time comparison assumes the manual operator is proficient. Differences in the operator's manual modeling speed could skew the speedup factor.

\section{Summary}
This chapter introduced a dataset-driven evaluation methodology to assess the proposed AI collaborator. By benchmarking 10 optimization tasks across three different LLM sizes, the evaluation provides quantitative insights into the system's efficiency gains (RQ1), architectural accuracy and safety (RQ2), and the effectiveness of the human-in-the-loop feedback mechanism for error recovery (RQ3).